% Préambule allégé utilisé UNIQUEMENT pour la version web (plasTeX).
% Le PDF continue d'utiliser preamble.tex.
\documentclass{book}
\usepackage{amsmath,amssymb,amsthm}
\usepackage{hyperref}
\usepackage{cleveref}
\usepackage{webfixes}

% --- Environnements (mêmes noms que preamble/box.tex) ---
\newtheorem{theorem}{Theorem}[section]
\newtheorem{lemma}[theorem]{Lemma}
\newtheorem{proposition}[theorem]{Proposition}
\newtheorem{corollary}[theorem]{Corollary}
\theoremstyle{definition}
\newtheorem{definition}[theorem]{Definition}
\newtheorem*{theorem*}{Theorem}
\newtheorem{customthm}{Theorem}
\newtheorem{customlemma}{Lemma}
\newtheorem{customproposition}{Proposition}

% --- Macros mathématiques (miroir de preamble/macros.tex) ---
\newcommand{\norm}[1]{\left\lVert#1\right\rVert}
\newcommand{\tpmod}[1]{\pmod{#1}}
\DeclareMathOperator{\ord}{ord}
\newcommand{\mathds}[1]{\mathbb{#1}}
\newcommand{\rcom}[1]{\textit{Rebecca: #1}}

% --- Commandes de mise en page du PDF, neutralisées pour le web ---
\newcommand{\intro}[1]{#1}
\newcommand{\startcontents}{}
\newcommand{\printcontents}[3]{}
\newcommand{\chaptertoc}{}
\newcommand{\mytitlepage}[1]{}
\newcommand{\upshape}{}
\newcommand{\index}[2][]{}
\newcommand{\todo}[1]{}
\newcommand{\lipsum}{}
\newcommand{\parencite}[1]{\cite{#1}}
\newcommand{\textcite}[1]{\cite{#1}}
\newenvironment{appendices}{\appendix}{}

% Correctifs pour plasTeX
